\section{Double slit and quantum eraser as objecthood budgeting}\label{sec:doubleslit}

The double-slit experiment is often framed as a ``wave versus particle'' dilemma.
In the SBT packaging language, a more precise description is: the double-slit experiment is a controlled demonstration of \emph{when a distinction becomes an object of the operative layer}.
The relevant distinction is ``which slit did the particle pass through?''
If this distinction is not packaged into a stable record, the layer's correct composition rule is coherent superposition (producing interference);
if the distinction is packaged, the correct rule is a classical mixture (no interference).

\subsection{A minimal model with explicit record overlap}

Let the slit degree of freedom be spanned by $\{|A\rangle,|B\rangle\}$, and let $|x\rangle$ denote a screen position basis.
After the slits, consider the coherent preparation
\[
|\psi\rangle = \tfrac{1}{\sqrt2}\left(|A\rangle + |B\rangle\right).
\]
Let $U$ be the propagation operator from the slit plane to the screen, and define the two path amplitudes as
\[
\psi_A(x) := \langle x|U|A\rangle,\qquad \psi_B(x) := \langle x|U|B\rangle.
\]
Without any which-path record, the detection probability at position~$x$ is
\[
P(x) = \left|\tfrac{1}{\sqrt2}(\psi_A(x)+\psi_B(x))\right|^2
     = \tfrac12\left(|\psi_A(x)|^2+|\psi_B(x)|^2\right)
       + \mathrm{Re}\!\left(\psi_A(x)^\ast\psi_B(x)\right).
\]
The last term is the interference term.

Now model any which-path record (including uncontrolled environmental coupling) by an environment register~$\mathcal{E}$ that becomes correlated with the path:
\[
|A\rangle|\mathcal{E}_0\rangle \mapsto |A\rangle|\mathcal{E}_A\rangle,\qquad
|B\rangle|\mathcal{E}_0\rangle \mapsto |B\rangle|\mathcal{E}_B\rangle.
\]
The joint state becomes
\[
|\Psi\rangle=\tfrac{1}{\sqrt2}\big(|A\rangle|\mathcal{E}_A\rangle + |B\rangle|\mathcal{E}_B\rangle\big).
\]
If the environment record is not retained or conditioned upon at the screen, the observed statistics are
\[
P(x)=\tfrac12\left(|\psi_A(x)|^2+|\psi_B(x)|^2\right)
+\mathrm{Re}\!\left(\psi_A(x)^\ast\psi_B(x)\,\langle \mathcal{E}_B|\mathcal{E}_A\rangle\right).
\]
Thus the interference visibility is controlled by the record overlap~$\gamma$:\cite{Englert1996Visibility}
\[
\gamma := \langle \mathcal{E}_B|\mathcal{E}_A\rangle.
\]
When $\gamma=1$, the record carries no distinguishability and interference is maximal; when $\gamma=0$, the record is perfectly distinguishing and interference vanishes.

\subsection{SBT interpretation: when a distinction becomes an object}

In SBT terms, ``which slit?'' is not automatically an object of the operative layer.
The distinction ``which slit?'' becomes an object precisely when the substrate supports a stable, distinguishable carrier of that distinction.
In the equation above, this occurs exactly in the limit $|\gamma|\to 0$:
the environment states become orthogonal, the distinction becomes recordable, and the interference term disappears.

Conversely, when $|\gamma|$ is near~$1$, the layer cannot stably split the description into ``path~$A$ occurred'' versus ``path~$B$ occurred'' without inventing surplus structure.
The correct layer description is then to compose the two alternatives coherently, producing interference.

\subsection{Reproducible visibility curves}

Figure~\ref{fig:ds-patterns} shows the screen patterns for several values of~$|\gamma|$ in a far-field minimal model.
Figure~\ref{fig:ds-visibility} shows that the extracted fringe visibility tracks~$|\gamma|$ as expected.
Both figures are generated deterministically by repository experiment \texttt{EXP-DS1} (see Appendix~\ref{app:repro}).

\begin{figure}[t]
  \centering
  \includegraphics[width=0.95\linewidth]{double_slit_patterns.png}
  \caption{Double slit patterns as the record overlap $|\gamma|$ decreases: interference fades exactly as a stable which-path record forms.}
  \label{fig:ds-patterns}
\end{figure}

\begin{figure}[t]
  \centering
  \includegraphics[width=0.75\linewidth]{double_slit_visibility.png}
  \caption{Extracted visibility versus record overlap magnitude $|\gamma|$ in the minimal model.}
  \label{fig:ds-visibility}
\end{figure}

\subsection{Quantum eraser as repackaging (not retrocausality)}

The same logic explains ``quantum eraser'' phenomena without invoking any backwards-in-time influence.\cite{ScullyDruhl1982,Kim2000DelayedChoice}
Marking which-path information correlates the path with an environment degree of freedom, driving~$|\gamma|$ toward~$0$ and eliminating interference in the \emph{unconditional} statistics.
However, if one subsequently conditions on an environment measurement performed in a basis that does \emph{not} preserve the which-path split (e.g., the $\{|\pm\rangle\}$ basis for an environment qubit),
then within each conditioned subensemble the effective overlap is restored and interference reappears, with complementary phase shifts between the two subensembles.

This is packaging language: different conditioning choices correspond to different ways of carving the joint state into record-level objects.
Nothing needs to ``change the past''; what changes is which distinctions are stabilized and reported at the layer.

\begin{figure}[t]
  \centering
  \includegraphics[width=0.95\linewidth]{quantum_eraser_unconditional.png}
  \caption{Quantum eraser: unconditional statistics after which-path marking (interference suppressed).}
  \label{fig:qe-uncond}
\end{figure}

\begin{figure}[t]
  \centering
  \includegraphics[width=0.95\linewidth]{quantum_eraser_conditional.png}
  \caption{Quantum eraser: conditional subensembles after measuring the environment in the $\{|\pm\rangle\}$ basis (interference restored with opposite phase).}
  \label{fig:qe-cond}
\end{figure}
